\chapter{Survey analysis}\label{ch:surveys}
\chapterstatus{written}{2026-11-14}

From mock or observed galaxy catalogues to clustering measurements: survey mocks from periodic boxes, the conversion of
sky coordinates to a Cartesian embedding, FKP power-spectrum multipoles with a local line of sight, the survey window and its
convolution of model spectra, and BAO reconstruction.

\begin{modulebox}{Survey}
\nsp{FML::SURVEY}: \file{Survey/GalaxiesToBox.h}, \file{Survey/SurveyMesh.h} (\cls{Geometry}, \cls{CatalogPoint}),
\file{Survey/SurveyPower.h} (\code{estimate_power}, \cls{PowerOptions}), \file{Survey/SurveyWindow.h} (\code{estimate_window},
\code{make_binned_window_matrix}), \file{Survey/SurveyReconstruction.h} (\code{reconstruct}, \cls{ReconstructionOptions});
\nsp{FML::SURVEY::MOCKS}: \file{Survey/SurveyMocks.h} (\code{make_survey}, \code{make_randoms}, \code{make_remapping}). The program
\file{Survey/make_survey}.\par
Documentation: \file{Survey/README.md}, \file{power_validation/}, \file{window_validation/}. Tests: \file{Survey/tests}; examples
\file{workflow_example}, \file{mock_example}, \file{mock_clustering_example}.
\end{modulebox}

\section{Survey mocks}\label{sec:survey-mocks}

\code{make_survey} turns objects in a periodic box (HOD galaxies, \cref{ch:galaxy-mocks}, or particles) into a survey seen from an observer:
the box is tiled around the observer, each object gets RA/Dec in a survey frame and the observed redshift
\begin{equation}
z_{\rm obs}=z_{\rm cos}+(1+z_{\rm cos})\frac{v_r}{c},\qquad\bm x_{\rm obs}=\chi(z_{\rm obs})\hat{\bm r},
\end{equation}
with $\chi(z)=\int c\,\dd z/H$ from a user-supplied $H(z)/c$, and is kept with probability (target $\bar n(z)$ / parent density) $\times$ footprint
completeness. Footprints are RA/Dec boxes, HEALPix maps \cite{Gorski2005} (the pixelisation is built in) or MANGLE polygons \cite{Swanson2008}.
Each object stores the selection probability, completeness and the expected density $\bar n$ of the final catalogue, for FKP weights later;
\code{make_randoms} makes randoms with the same selection. Random decisions use Philox streams keyed by object and replica, so the mock does not
depend on the number of tasks. Repetition of structures when $\chi(z_{\max})>L/2$ is avoided with the volume-preserving cuboid remapping of Carlson
and White \cite{CarlsonWhite2010} (\code{make_remapping} with an integer unimodular matrix). Evolving lightcones from \file{SimulationCode} are
used directly (\code{make_survey_lightcone}). The program \file{make_survey} does all this from a Lua parameter file. \file{GalaxiesToBox.h}
converts (RA, Dec, $z$) catalogues to Cartesian positions.

\section{Power-spectrum multipoles}\label{sec:survey-power}

With data and randoms (positions in physical units inside a padded FFT box, effective weights and $\bar n$), the FKP field \cite{Feldman1994} and
its normalisation are
\begin{equation}
F(\bm x)=n_w^D(\bm x)-\alpha n_w^R(\bm x),\qquad\alpha=\frac{\sum_Dw}{\sum_Rw},\qquad I_2=\alpha\sum_R\bar n\,w^2,
\end{equation}
with FKP weights $w=1/(1+\bar nP_0)$ supplied by the user. The multipoles with a local first-point line of sight \cite{Yamamoto2006,Bianchi2015,
Scoccimarro2015} are
\begin{equation}
A_\ell(\bm k)=\sum_{\bm x}F(\bm x)\mathcal L_\ell(\hat{\bm k}\cdot\hat{\bm x})e^{-i\bm k\cdot\bm x},\qquad
\hat P_\ell(k)=\frac{2\ell+1}{I_2}\left\langle{\rm Re}\left[A_0(\bm k)A_\ell^*(\bm k)\right]\right\rangle_{\rm shell}-N_\ell,
\end{equation}
evaluated with the tensor (Cartesian-moment) FFTs of $\hat x_i\hat x_j\cdots F$; a fixed line of sight is an option. $\ell=0,2,4$. The weighted Poisson
noise $N_\ell$ is computed from the catalogue's line-of-sight moments (data weights squared plus $\alpha^2$ random weights squared). CIC compensation
divides by the window, and the shot noise includes the aliased images of the CIC-painted noise \cite{Jing2005}, $\prod_i[1-\frac23\sin^2(k_ih/2)]$
(without it 5\% of the shot noise remains at $kh=1.6$--2.6); there is no interlacing, so scales well below the Nyquist frequency should be used.

\section{The window and convolved models}\label{sec:survey-window}

The window multipoles of the randoms' selection $W=\alpha n_w^R/V_{\rm cell}$ are
\begin{equation}
Q_L(s)=\frac{2L+1}{I_2}\left\langle\int\dd^3x\,W(\bm x)W(\bm x+\bm s)\mathcal L_L(\hat{\bm s}\cdot\hat{\bm x})\right\rangle_{\hat{\bm s}},
\end{equation}
up to $L=8$, from FFT cross-correlations of the line-of-sight tensors, with the CIC self-pairs of each random removed exactly. Because the lattice vectors
of a separation shell are not isotropic, each shell uses $\langle GF\rangle-\langle G\rangle[\langle F\rangle-\langle F\rangle_{\rm sphere}]$, exact for a pair
function isotropic on the shell (this reduced the leakage of the monopole into $Q_4$ from $\sim4$ to $\sim0.01$ in a test). The convolved model multipoles
follow Wilson et al.\ \cite{Wilson2017}: model $P_\ell$ bandpowers are transformed to $\xi_\ell$, multiplied by $Q_L$ with the Legendre triple-product
coefficients, and transformed back with spherical Bessel kernels; \code{make_binned_window_matrix} averages the output over the measured shells, giving a
matrix from input to measured bandpowers. Input bins must extend well beyond the last measured bin. An optional global integral constraint makes the
convolved monopole vanish at $k=0$. Wide-angle terms beyond the local plane-parallel convention are not included.

\section{BAO reconstruction}\label{sec:survey-reconstruction}

Reconstruction \cite{Eisenstein2007,Padmanabhan2012} solves for the Zel'dovich displacement of the smoothed tracer field
$\delta_t=D_s/(\alpha R_s)-1$ (Gaussian smoothing, cells with too few randoms or outside a mask set to zero),
\begin{equation}
-\nabla\cdot\left[\left(\mathsf I+\frac fb\hat{\bm n}\hat{\bm n}\right)\nabla\phi\right]=\frac{\delta_t}{b},\qquad\bm\Psi=\nabla\phi,
\end{equation}
with the local radial or a fixed line of sight, by preconditioned conjugate gradients (a positive-definite finite-volume stencil, a Fourier Laplacian
preconditioner). Data move by $-\bm\Psi-f(\bm\Psi\cdot\hat{\bm n})\hat{\bm n}$; randoms by the same (RecSym) or by $-\bm\Psi$ (RecIso). The shifted
catalogues go back into the power estimator. A solver report gives the convergence and the removed source mean; non-convergence throws.

\begin{cppcode}[Survey multipoles before and after reconstruction]
#include <FML/Survey/SurveyPower.h>
#include <FML/Survey/SurveyReconstruction.h>
using namespace FML::SURVEY;
Geometry geo;  geo.nmesh = 256;  geo.boxsize = 2000;  geo.minimum = {-1000, -1000, -1000};
PowerOptions opt;  opt.edges = edges;  opt.ells = {0, 2, 4};
auto pre = estimate_power(data, randoms, geo, opt);
ReconstructionOptions rec;  rec.bias = 2.0;  rec.growth_rate = 0.8;  rec.smoothing_radius = 15;
auto shifted = reconstruct(data, randoms, geo, rec);
auto post = estimate_power(shifted.data, randoms, geo, opt, &shifted.randoms);
\end{cppcode}

\begin{validation}[Validation: surveys]
FFT multipoles against a direct DFT; raw fixed and local line-of-sight multipoles against pypower with matched conventions to $10^{-15}$ of the monopole; mesh
convergence against a catalogue DFT (0.02--0.07\% at $128^3$ for the test catalogue); window multipoles against explicit weighted random pairs and the window
matrix against direct angular/Bessel integration and pypower; displacement solves against manufactured solutions (an independent spectral solution agreed to a
correlation of 0.99995); RecIso/RecSym, masks, MPI ownership. Mocks: $\chi(z)$ against EdS to $4\times10^{-9}$, coordinate round trips, the footprint area and
$n(z)$ of data and randoms, the remapping and HEALPix pixels against healpy. End to end (\file{mock_clustering_example}): a Zel'dovich box to a half-sky mock gives
$P_0$ within a few per cent of the window-convolved Kaiser model for $0.015<k<0.1\hMpc$ (mean of 3 realisations).
\end{validation}
